\documentclass[11pt]{article}
\usepackage[margin=1in]{geometry}
\usepackage{amsmath,amssymb,amsthm}
\usepackage{hyperref}
\newtheorem{theorem}{Theorem}
\newtheorem{lemma}[theorem]{Lemma}
\newtheorem{proposition}[theorem]{Proposition}
\newtheorem{corollary}[theorem]{Corollary}
\theoremstyle{definition}
\newtheorem{definition}[theorem]{Definition}
\theoremstyle{remark}
\newtheorem{remark}[theorem]{Remark}
\newcommand{\Z}{\mathbb{Z}}
\newcommand{\Jac}{\operatorname{Jac}}
\newcommand{\Pic}{\operatorname{Pic}}

\title{Disproof of Conjecture 00000001230:\\
the Jacobian of the $2\times2$ grid is $\Z/4$, which is not a sum of $\Z/\gcd(S)$}
\author{jilint777}
\date{2026-10-04}

\begin{document}
\sloppy
\maketitle

\begin{abstract}
Conjecture 00000001230 asserts that the Jacobian (sandpile group) of the grid
graph $[m_1]\times\dots\times[m_n]$ is isomorphic to a group
$\bigoplus_S(\Z/\gcd(S))^{e(S)}$, where $S$ runs over the subsets of
$\{m_1,\dots,m_n\}$. The exponents $e(S)$ are not specified. We show that no
choice of exponents works. The $2\times2$ grid is the $4$-cycle $C_4$, and its
Jacobian $\Z/4$ has an element of order $4$. For $(m_1,m_2)=(2,2)$, every
modulus $\gcd(S)$ is $0$ or $2$. In any product of copies of $\Z$ and $\Z/2$,
an element $x$ with $4x=0$ also has $2x=0$. So $\Z/4$ does not even embed into
any such group. The same graph defeats the edge-count convention
$[m]=$ path with $m$ edges, where the moduli are $0$ and $1$. The cube
$[2]^3$ (Jacobian $\Z/2\oplus\Z/8\oplus\Z/24$) gives a counterexample in
dimension $3$. The disproof is formalized in Lean~4 (core library only), with
the grid graph, its Laplacian, the sandpile group and the conjectured groups
all defined from scratch.
\end{abstract}

\section{The conjecture and its precise reading}
\begin{quote}
\emph{Definition:} $\Jac(K_{m,n})$ is the sandpile group (Jacobian) of the
complete bipartite graph. \emph{Conjecture:} The Jacobian of the
$n$-dimensional grid graph $[m_1]\times\dots\times[m_n]$ is isomorphic to the
explicit decomposition $\bigoplus_S(\Z/\gcd(S))^{e(S)}$ over subsets $S$ of
the $m_i$; the two-dimensional case is given explicitly by the $K_{m,n}$
formula and torus gluing.
\end{quote}
The Chinese version says the same. The statement never defines $e(S)$, so we
read it in the weakest possible form:
\begin{quote}
\textbf{Claim.} For every $n\ge1$ and every $m_1,\dots,m_n\ge1$ there exist
exponents $e(S)\ge0$, one for each subset $S$, such that
$\Jac([m_1]\times\dots\times[m_n])\cong\bigoplus_S(\Z/\gcd(S))^{e(S)}$.
\end{quote}
Any explicit formula for $e(S)$, including whatever the second clause
(``given explicitly by the $K_{m,n}$ formula and torus gluing'') is meant to
produce in dimension $2$, is a special case. Refuting the Claim therefore
refutes the conjecture, whatever the exponents are meant to be.
The second clause is too vague to state precisely. We do not need it, since
the conjecture is a conjunction and the first clause already fails.

\paragraph{Conventions covered.} The counterexample defeats all of the
following readings at the same time (Section~\ref{sec:readings}):
\begin{itemize}
\item $[m]$ is the path with $m$ \emph{vertices} (standard), or the path with
$m$ \emph{edges};
\item $\gcd(\emptyset)=0$ (so $\Z/\gcd(\emptyset)=\Z$), $\gcd(\emptyset)=1$, or
the empty set is skipped;
\item ``subsets of the $m_i$'' are subsets of the index set $\{1,\dots,n\}$
(so repeated values give repeated summands) or of the value set
$\{m_1,\dots,m_n\}$;
\item the exponents are finite, or arbitrary cardinals (infinite direct sums);
\item ``Jacobian'' is the sandpile group $\Z^{V\setminus\{s\}}/\tilde\Delta\Z^{V\setminus\{s\}}$,
the Baker--Norine Jacobian $\Pic^0=\mathrm{Div}^0/\mathrm{Prin}$, or even the
full cokernel $\Z^V/\Delta\Z^V$;
\item dimension $n=2$, or the claim restricted to $n\ge3$;
\item ``$\gcd$'' read as a typo for $\operatorname{lcm}$ or product. Here the
$2\times2$ grid is not enough, but the $2\times3$ grid is
(Section~\ref{sec:readings}; also in Lean).
\end{itemize}

\section{Definitions}
\begin{definition}
For $m_1,\dots,m_n\ge1$, the grid graph $[m_1]\times\dots\times[m_n]$ has
vertex set $\{(x_1,\dots,x_n):0\le x_i<m_i\}$. Two vertices are adjacent iff
their $\ell^1$ distance is $1$, that is, they differ by $1$ in exactly one
coordinate. Equivalently, it is the Cartesian product $P_{m_1}\square\cdots\square P_{m_n}$
of paths.
\end{definition}

\begin{definition}
Let $G=(V,E)$ be a connected graph with Laplacian $\Delta=D-A$, and fix a sink
$s\in V$. The \emph{reduced Laplacian} $\tilde\Delta$ is $\Delta$ with the row
and column of $s$ deleted. The \emph{sandpile group} (Jacobian, critical group)
is
\[
K(G)=\Z^{V\setminus\{s\}}/\tilde\Delta\,\Z^{V\setminus\{s\}} .
\]
Its order is $\det\tilde\Delta$, the number of spanning trees
(matrix-tree theorem). Up to isomorphism it does not depend on $s$. It is
isomorphic to the Baker--Norine Jacobian
$\Pic^0(G)=\mathrm{Div}^0(G)/\Delta\Z^V$, where $\mathrm{Div}^0$ is the set of
integer vectors with coordinate sum $0$. Also $\Pic(G)=\Z^V/\Delta\Z^V\cong\Z\oplus K(G)$.
\end{definition}
Throughout, $\Z/0=\Z$ and $\Z/1=0$.

\section{The obstruction}
\begin{lemma}\label{lem:target}
Let $(d_i)_{i\in I}$ be any family, indexed by any set $I$, with
$d_i\in\{0,1,2\}$. Let $P=\prod_{i\in I}\Z/d_i$. If $x\in P$ and $4x=0$, then
$2x=0$. The same holds in every subgroup of $P$, in particular in the direct
sum $\bigoplus_{i\in I}\Z/d_i$.
\end{lemma}
\begin{proof}
Work coordinatewise. If $d_i=0$, then $4x_i=0$ in $\Z$ forces $x_i=0$. If
$d_i=2$ or $d_i=1$, then $2x_i=0$ always.
\end{proof}

\begin{lemma}\label{lem:transfer}
Let $G$ be an abelian group with an element $g$ such that $4g=0$ and $2g\ne0$.
Then for every family as in Lemma~\ref{lem:target} there is no injective
homomorphism $G\to\prod_i\Z/d_i$. In particular $G\not\cong\bigoplus_i\Z/d_i$.
\end{lemma}
\begin{proof}
Let $\varphi$ be an injective homomorphism. Then $4\varphi(g)=\varphi(4g)=0$, so
$2\varphi(g)=0$ by Lemma~\ref{lem:target}. So $\varphi(2g)=0=\varphi(0)$, and
injectivity gives $2g=0$, a contradiction.
\end{proof}
In Lean, a homomorphism is just a map $\varphi$ with
$\varphi(a+b)=\varphi(a)+\varphi(b)$, and $\varphi(0)=0$ is derived from
this. Lemma~\ref{lem:target} also gives cancellation, $y+y=y\Rightarrow y=0$,
in $\prod_i\Z/d_i$.

\section{\texorpdfstring{The $2\times2$ grid}{The 2x2 grid}}
The grid $[2]\times[2]$ has vertices $(0,0),(0,1),(1,0),(1,1)$ and edges
$(0,0)(0,1)$, $(0,0)(1,0)$, $(0,1)(1,1)$, $(1,0)(1,1)$. It is the $4$-cycle
$C_4$. Its Laplacian, in this vertex order, is
\[
\Delta=\begin{pmatrix}2&-1&-1&0\\-1&2&0&-1\\-1&0&2&-1\\0&-1&-1&2\end{pmatrix}.
\]
With sink $s=(0,0)$, the reduced Laplacian is
\[
\tilde\Delta=\begin{pmatrix}2&0&-1\\0&2&-1\\-1&-1&2\end{pmatrix},\qquad
\det\tilde\Delta=4 .
\]

\begin{proposition}\label{prop:c4}
Let $g=[e_0]\in K(C_4)$ be the class of one chip on $(0,1)$. Then $4g=0$ and
$2g\neq0$. In fact $K(C_4)\cong\Z/4$, via $x\mapsto 3x_0+x_1+2x_2\bmod 4$.
\end{proposition}
\begin{proof}
$\tilde\Delta(3,1,2)^{\mathsf T}=(6-2,\;2-2,\;-3-1+4)^{\mathsf T}=(4,0,0)^{\mathsf T}=4e_0$, so $4g=0$.
Suppose $\tilde\Delta z=2e_0$ with $z\in\Z^3$, that is,
$2z_0-z_2=2$, $2z_1-z_2=0$ and $-z_0-z_1+2z_2=0$. The second equation makes $z_2$
even, say $z_2=2a$. Then $z_1=a$ and $z_0=1+a$, and the third equation becomes
$2a=1$, which is impossible. So $2g\ne0$.

For the isomorphism, the vector $(3,1,2)$ is the first row of
$\operatorname{adj}\tilde\Delta=4\tilde\Delta^{-1}$. It annihilates every column
of $\tilde\Delta$ modulo $4$ (columns $(2,0,-1)$, $(0,2,-1)$, $(-1,-1,2)$ give
$4$, $0$, $0$). So $\psi(x)=3x_0+x_1+2x_2\bmod4$ is well defined on $K(C_4)$.
It is onto, since $\psi(3a,0,0)=9a\equiv a$. It is injective: if
$w=x-y$ has $\psi(w)\equiv0$, write $3w_0+w_1+2w_2=4q$. Then $w_0-w_1$ is even,
say $2r$, and $z=(q,\;q-r,\;r+w_1+w_2)$ is an integer solution of
$\tilde\Delta z=w$.
\end{proof}

\begin{theorem}\label{thm:main}
For $(m_1,m_2)=(2,2)$, no choice of exponents gives
$K([2]\times[2])\cong\bigoplus_S(\Z/\gcd S)^{e(S)}$. There is not even an
injective homomorphism from $K([2]\times[2])$ into such a group. Hence the
Claim, and with it Conjecture 00000001230, is false.
\end{theorem}
\begin{proof}
The subsets of $\{m_1,m_2\}=\{2,2\}$ (as index subsets) are $\emptyset$,
$\{m_1\}$, $\{m_2\}$ and $\{m_1,m_2\}$, with gcds $0,2,2,2$. So every modulus
is $0$ or $2$. Apply Lemma~\ref{lem:transfer} to $G=K(C_4)$ with the element of
Proposition~\ref{prop:c4}.
\end{proof}

\section{Robustness}\label{sec:readings}
\paragraph{Edge-count convention.} If $[m]$ is the path with $m$ edges, the
$4$-cycle is $[1]\times[1]$. The moduli are then $\gcd(\emptyset)=0$ and
$\gcd\{1\}=\gcd\{1,1\}=1$. All of them lie in $\{0,1,2\}$, so Lemma~\ref{lem:transfer}
applies again. (A product of copies of $\Z$ and $0$ is torsion-free.)

\paragraph{Empty set and value-set reading.} With $\gcd(\emptyset)=1$, the
moduli are $\{1,2\}$. If the empty set is skipped, or $S$ ranges over subsets
of the value set $\{2\}$, the available moduli form a subset of $\{0,2\}$.
Every case is covered by Lemma~\ref{lem:transfer}, which allows an arbitrary
family of moduli from $\{0,1,2\}$, including the choice $e(S)=0$.

\paragraph{Infinite exponents.} Lemma~\ref{lem:transfer} allows any index set $I$.
In Lean the index type $\iota$ is arbitrary.

\paragraph{Other meanings of ``Jacobian''.} In $\Pic(C_4)=\Z^V/\Delta\Z^V$, the
degree-zero class $p=[(0,1)-(0,0)]$ satisfies
$\Delta(0,3,1,2)^{\mathsf T}=(-4,4,0,0)^{\mathsf T}=4\,((0,1)-(0,0))$, so $4p=0$. Moreover
$2p\ne0$: if $\Delta z=2((0,1)-(0,0))$, then subtracting the constant $z_{(0,0)}$
from $z$ (the all-ones vector lies in $\ker\Delta$) and reading the last three
rows gives $\tilde\Delta z'=2e_0$, which is impossible. So
Lemma~\ref{lem:transfer} applies both to $\Pic(C_4)$ and to its subgroup
$\Pic^0(C_4)$, which contains $p$. Lean formalizes both.

\paragraph{Dimension three.} If the claim is meant only for $n\ge3$, take the
cube $[2]\times[2]\times[2]$. Its moduli are again $0$ and $2$ (respectively $0$ and $1$
for the edge convention $[1]^3$). Order the non-sink vertices
$(0,0,1),(0,1,0),(0,1,1),(1,0,0),(1,0,1),(1,1,0),(1,1,1)$, with sink $(0,0,0)$.
Then $\tilde\Delta(3,3,6,2,3,3,4)^{\mathsf T}=8e_2$, where $e_2$ is the chip at
$(0,1,1)$. The unique rational solution of $\tilde\Delta z=4e_2$ is not
integral ($\det\tilde\Delta=384$; the solution is $\frac12(3,3,6,2,3,3,4)$).
So $g=[2e_2]$ has $4g=0\ne2g$, and Lemma~\ref{lem:transfer} applies. In fact
$K([2]^3)\cong\Z/2\oplus\Z/8\oplus\Z/24$.

\paragraph{The failure is generic.} By Smith normal form (\texttt{verify.py}),
\begin{center}
\begin{tabular}{lll}
grid & $K$ & moduli (vertex / edge convention)\\\hline
$2\times2$ & $\Z/4$ & $\{0,2\}$ / $\{0,1\}$\\
$2\times3$ & $\Z/15$ & $\{0,1,2,3\}$ / $\{0,1,2\}$\\
$3\times3$ & $\Z/8\oplus\Z/24$ & $\{0,3\}$ / $\{0,2\}$\\
$2\times4$ & $\Z/56$ & $\{0,2,4\}$ / $\{0,1,3\}$\\
$3\times4$ & $\Z/2415$ & $\{0,1,3,4\}$ / $\{0,1,2,3\}$\\
$2\times2\times2$ & $\Z/2\oplus\Z/8\oplus\Z/24$ & $\{0,2\}$ / $\{0,1\}$
\end{tabular}
\end{center}
A finite group $\bigoplus_d(\Z/d)^{e_d}$ has no $\Z$ summand. Its exponent
divides the lcm of the nonzero moduli. In every row, the largest invariant factor of $K$
does not divide that lcm. The order of $K$ is the number of spanning trees,
and it is typically divisible by primes that do not divide any $m_i$, such as $5$
in $\Z/15$ or $7$ in $\Z/2415$. Such primes can never come from the moduli.

\paragraph{``gcd'' read as lcm or product.} If $\gcd(S)$ were a typo for
$\operatorname{lcm}(S)$ or $\prod S$, the $2\times2$ grid would \emph{not} be a
counterexample, because $\Z/(2\cdot2)=\Z/4$. The $2\times3$ grid is one. Its
Jacobian is $K([2]\times[3])\cong\Z/15$. With sink $(0,0)$ and non-sink
vertices ordered $(0,1),(0,2),(1,0),(1,1),(1,2)$, we have
$\tilde\Delta(3,5,2,4,7)^{\mathsf T}=5e_4$, where $e_4$ is the chip at $(1,2)$.
Since $\det\tilde\Delta=15$, the unique rational solution of $\tilde\Delta z=e_4$ is
$\frac15(3,5,2,4,7)$, which is not integral. So $g=[e_4]$ has order exactly $5$. Now consider the
gcds, lcms and products of all subsets of $\{2,3\}$, or of $\{1,2\}$ in the
edge convention. With $\gcd\emptyset=0$ and $\operatorname{lcm}\emptyset=\prod\emptyset=1$,
each of these is $0$ or a divisor of $6$. This also covers ``$K_{m,n}$-style'' moduli
$m,n,mn=2,3,6$. In $\prod_i\Z/d_i$ with every $d_i=0$ or $d_i\mid6$, $5x=0$
forces $x=0$. The proof of Lemma~\ref{lem:transfer}, with $5$ in place of
$4$ and $2$, then shows that $K([2]\times[3])$ embeds in no such group. This is
formalized in Lean (\texttt{no\_iso\_23\_gcd\_lcm\_prod},
\texttt{no\_iso\_12\_edges\_gcd\_lcm\_prod}) and checked in \texttt{verify.py}.

\begin{remark}
The conjecture mentions ``torus gluing''. Its object is the grid graph, not
the torus, but the torus readings fail too. The simple graph $C_2\square C_2$
is literally $C_4=[2]\times[2]$, so it is the same counterexample. The
multigraph $C_2\square C_2$ (each edge doubled) has
$K=\Z/2\oplus\Z/2\oplus\Z/8$, which contains an element of order $8$, against
moduli $\{0,2\}$. The multigraph $C_2\square C_3$ has $K=\Z/7\oplus\Z/42$, against
moduli $\{0,1,2,3\}$. Also $K(C_3\square C_3)=(\Z/6)^2\oplus(\Z/18)^2$ against
moduli $\{0,3\}$, and $K(C_3\square C_4)=(\Z/5)^2\oplus\Z/35\oplus\Z/420$ against
$\{0,1,3,4\}$. The strong product $P_2\boxtimes P_2=K_4$ has
$K=(\Z/4)^2$, which fails against moduli $\{0,2\}$ for the same order-$4$ reason.
These facts are checked in \texttt{verify.py} only, not in Lean.
\end{remark}

\section{Lean formalization}
The directory \texttt{lean4/} is a self-contained Lean~4.19.0 project (core
library only, no Mathlib).
\begin{itemize}\raggedright
\item \texttt{gridVerts ms} lists the vertices of the grid for an
\emph{arbitrary} list \texttt{ms} $=[m_1,\dots,m_n]$. \texttt{adj} is
adjacency at $\ell^1$ distance $1$, \texttt{deg} the degree, \texttt{lap} the
Laplacian, and \texttt{redLap} the reduced Laplacian with sink $(0,\dots,0)$.
All of these are general definitions. The concrete matrices are obtained by
evaluating them (\texttt{redLap\_22}, \texttt{lap\_22} and \texttt{redLap\_222},
by \texttt{decide}).
\item \texttt{Coker k M} is the quotient type $\Z^k/M\Z^k$, where $x\sim y$
iff $x-y=Mz$ for some $z\in\Z^k$. Lean proves that this is an equivalence
relation compatible with addition. \texttt{Sandpile ms} is
\texttt{Coker} $(|V|-1)$ \texttt{(redLap ms)}, and \texttt{Pic ms} is
\texttt{Coker} $|V|$ \texttt{(lap ms)}. \texttt{Pic0} is the subtype of
$\Pic([2]\times[2])$ of degree-zero classes; the degree is well defined because
the columns of $\Delta$ sum to $0$.
\item \texttt{ZM d} is $\Z/d$, represented by the canonical representatives
$a$ with $a\bmod d=a$ (so \texttt{ZM 0} $=\Z$). \texttt{subsetsL} lists the
$2^n$ sub-lists, and \texttt{gcdL} is the gcd of a list with
$\gcd(\emptyset)=0$. \texttt{Target ms e} is the product over pairs
(subset $k$, copy $c<e(k)$) of $\Z/\gcd(S_k)$, that is,
$\bigoplus_S(\Z/\gcd S)^{e(S)}$. \texttt{IsIso} means additive and bijective.
\item \texttt{ClaimV} (vertex convention), \texttt{ClaimE} (edge convention)
and \texttt{ClaimV3} (vertex convention, $n\ge3$) state the Claim. The main
theorems are:
\begin{center}
\begin{tabular}{l}
\texttt{conjecture\_00000001230\_false : $\neg$ ClaimV}\\
\texttt{conjecture\_00000001230\_false\_edges : $\neg$ ClaimE}\\
\texttt{conjecture\_00000001230\_false\_dim3 : $\neg$ ClaimV3}
\end{tabular}
\end{center}
\item The explicit instances are \texttt{no\_iso\_22}, \texttt{no\_iso\_111\_edges},
\texttt{no\_iso\_22\_Pic0} and \texttt{no\_iso\_22\_Pic}. The strong forms
\texttt{no\_injective\_hom\_C4}, \texttt{\_Q3}, \texttt{\_PicC4} and
\texttt{\_Pic0C4} exclude injective additive maps into
$\prod_{i\in\iota}\Z/d_i$ for an arbitrary type $\iota$ and arbitrary
$d_i\in\{0,1,2\}$.
\item Witnesses: \texttt{g22\_four}, \texttt{g22\_two}, \texttt{p22\_four},
\texttt{p22\_two}, \texttt{g222\_four} and \texttt{g222\_two}. Each $4g=0$ is
proved with the explicit vector $z$. Each $2g\ne0$ is proved by
\texttt{omega} on the linear system.
\item Lcm/product reading: \texttt{redLap\_23} (the $5\times5$ reduced Laplacian of
$[2]\times[3]$, by \texttt{decide}), \texttt{g23\_five} ($5e_4=\tilde\Delta(3,5,2,4,7)$) and
\texttt{g23\_ne} ($e_4\notin\operatorname{im}\tilde\Delta$, by \texttt{omega}).
\texttt{no\_injective\_hom\_23} allows any moduli in $\{0,1,2,3,6\}$.
\texttt{no\_iso\_23\_gcd\_lcm\_prod} and \texttt{no\_iso\_12\_edges\_gcd\_lcm\_prod}
refute $K\cong\bigoplus_S(\Z/f(S))^{e(S)}$ for $f\in\{\gcd,\operatorname{lcm},\prod\}$
and every $e$, in both conventions.
\item Non-vacuity and sanity: \texttt{sandpile\_C4\_iso\_Z4} constructs the
isomorphism $K([2]\times[2])\cong\Z/4$ of Proposition~\ref{prop:c4}. This
shows that \texttt{Sandpile} is the right group and that \texttt{IsIso} can be
satisfied. \texttt{sandpile\_P3\_trivial} shows that the Claim's predicate
holds for the path $[3]$ (a tree; $e\equiv0$). So the Claim is not refuted
for a trivial reason.
\end{itemize}
\texttt{lake build} succeeds with no errors or warnings. There is no
\texttt{sorry}, no \texttt{native\_decide}, and no added axioms.
\texttt{\#print axioms} reports only \texttt{propext} and \texttt{Quot.sound}.

\paragraph{Limitations.} The second clause (the $2$-dimensional formula via
$K_{m,n}$ and torus gluing) is not formalized, because it is not stated
precisely. It is not needed, since the first clause already fails. The value-set reading of
``subsets'' and the convention $\gcd(\emptyset)=1$ are not separate Lean
predicates. They are covered by the strong lemmas, which allow every family of
moduli in $\{0,1,2\}$. The isomorphisms $K\cong\Pic^0$ and $\Pic\cong\Z\oplus K$
are not needed; Lean treats the three groups separately. The torus and strong-product
remark and the table above are checked only in \texttt{verify.py}.

\section{Independent check}
\texttt{verify.py} (Python 3 standard library) builds the grids from
coordinates. It computes Smith normal forms of the reduced and full Laplacians
with exact integer arithmetic and recomputes the witnesses. For each listed grid and both
conventions, it checks that the largest invariant factor of $K$ does not
divide the lcm of the nonzero moduli. It checks the lcm/product reading on the $2\times3$ grid. As sanity checks, it
verifies that paths (trees) have trivial Jacobian, so the predicate holds for
them, and it checks the torus and strong-product remark. It also notes that the
$2\times2$ grid alone would not refute the product reading.

\section{Reproduction}
\begin{verbatim}
cd lean4 && lake build     # Lean 4.19.0, core only
python3 verify.py          # Python 3 standard library
pdflatex report.tex
\end{verbatim}

\begin{thebibliography}{9}
\bibitem{bn} M.~Baker and S.~Norine, Riemann--Roch and Abel--Jacobi theory on a
finite graph, \emph{Adv. Math.} 215 (2007), 766--788.
\bibitem{kl} C.~J.~Klivans, \emph{The Mathematics of Chip-Firing}, CRC Press, 2018.
\bibitem{lor} D.~Lorenzini, Smith normal form and Laplacians, \emph{J. Combin.
Theory Ser. B} 98 (2008), 1271--1300.
\end{thebibliography}
\end{document}
